% !TEX root = ../main.tex

%* 设置页面布局
\geometry{
    %? 补药乱加空行
    a4paper,            % 纸张大小
    %? 双面打印
    inner=2.5cm,        % 内边距（靠装订线的边距）
    outer=2.5cm,        % 外边距
    %? 单面打印，少用
    % left=3cm,         % 左边距（靠装订线的边距）
    % right=2.5cm,      % 右边距
    top=3cm,            % 上边距
    bottom=2.5cm,       % 下边距
    headsep=1cm,        % 页眉与正文之间的距离
    headheight=14pt,    % 页眉高度
    footskip=1.5cm,     % 页脚与正文之间的距离
    bindingoffset=0cm   % 装订线偏移
}

%* 设置标题格式
\ctexset{
    section = {
        format = \raggedright\zihao{-3},
        number = \chinese{section},
        aftername = \hspace{1em}
    },
    subsection = {
        format = \raggedright\zihao{4},
        number = \arabic{subsection},
        aftername = \hspace{0.8em}
    },
    subsubsection = {
        format = \raggedright\zihao{-4},
        number = \arabic{section}.\arabic{subsection}.\arabic{subsubsection},
        aftername = \hspace{0.5em}
    }
}

%* 设置行距
\setstretch{1.25}       % 设置全局行距为1.25倍

%* 设置页眉页脚 
\pagestyle{fancy}                       % 设置页眉页脚样式
% \renewcommand{\headrulewidth}{0pt}    % 设置页眉线宽
\fancyhf{}                              % 清除默认设置
\fancyhead[R]{}                         % 页眉右侧
\fancyfoot[C]{\thepage}                 % 当前页码

%* 设置页码样式
% \pagenumbering{gobble}    % 关闭
% \pagenumbering{roman}
\pagenumbering{arabic}

%* 设置目录深度
\setcounter{tocdepth}{3}

%* 设置参考文献标题格式
% \renewcommand{\refname}{\zihao{5}\bfseries\heiti 参考文献：}

%* 设置多栏格式
% \setlength{\columnsep}{0.8cm}   % 设置两栏之间的间距
% \columnseprule 0.4pt          % 如果需要两栏中间有分隔线（可选，默认为0）